\documentclass[11pt]{article}
\usepackage{fontspec}
\usepackage[french,english]{babel}
\usepackage{amsmath}
\usepackage{amssymb}
\usepackage{amsthm}
\usepackage{mathrsfs}
\usepackage{tikz-cd}
\usepackage{enumitem}

\frenchsetup{AutoSpacePunctuation=false}
\tikzcdset{ampersand replacement=\&}

\newcommand{\printedpage}[1]{\marginpar{\footnotesize p.~#1}}
\newcommand{\msp}[1]{\marginpar{\footnotesize [#1]}}
\newcommand{\Ho}{\mathrm{Ho}}
\newcommand{\Fib}{\mathrm{Fib}}

\begin{document}

\printedpage{165}
% head: A. Grothendieck : Dérivateurs, Ch. XII, Édition des Universités de Montpellier II et Paris VII
% foot: Version 30 mars 2003

\textit{The functorial law is obtained with the help of lemma 1 p. 249, which allows us to restrict ourselves to the arguments $\begin{tikzcd}[ampersand replacement=\&] X_1 \arrow[d, "p_1"'] \& \\ S \& X_2 \arrow[l, "p_2"] \end{tikzcd}$ such that $p_1,p_2\in \mathit{Fib}_W$, so that we find a functor}
\[
\begin{tikzcd}[ampersand replacement=\&]
X_1 \arrow[d] \& \\
S \& X_2 \arrow[l]
\end{tikzcd}
\qquad\longmapsto\qquad
\begin{tikzcd}[ampersand replacement=\&]
X_1 \arrow[d] \& X_1\times_S X_2 \arrow[l] \arrow[d] \\
S \& X_2, \arrow[l]
\end{tikzcd}
\]
\textit{compatible with the localizers, and passing to the categories of fractions} $\Ho_W(\Psi)$ \textit{and} $\Ho_W(\tilde{\Psi})$ \textit{respectively.}

\textbf{NB} By proceeding in a similar, simpler way (with no need to resort to lemma 1), we find the existence of a \textit{left adjoint} $i_!$ to $i^* : \Ho_W(\Psi)\longrightarrow \Ho_W(\tilde{\Psi})$, subject to the existence of an initial object $\varphi$ in $\mathcal{M}$. This functor $i_!$ is given by % [sic] Cet foncteur -> Ce foncteur, "foncteur" begins with a consonant

\msp{page 255}
\[
i_!\left(
\begin{tikzcd}[ampersand replacement=\&]
X_1 \arrow[d] \& \\
S \& X_2 \arrow[l]
\end{tikzcd}
\right)
\qquad = \qquad
\left(
\begin{tikzcd}[ampersand replacement=\&]
X_1 \arrow[d] \& \varphi \arrow[l] \arrow[d] \\
S \& X_2 \arrow[l]
\end{tikzcd}
\right),
\]
the functorial law being evident.

This is not yet the end of the sorites! We must still interpret the composite functor
\[
\Ho_W(\Psi)\xrightarrow{\ i_*\ }\Ho_W(\tilde{\Psi})\xrightarrow{\ \varepsilon_3^*\ }\Ho_W(e)=\Ho_W,
\]
where
\[
\varepsilon_3 : e\xrightarrow{\ \sim\ }\{3\}
\qquad\longrightarrow\qquad
\begin{tikzcd}[ampersand replacement=\&, column sep=small, row sep=small]
\& 1 \arrow[dl] \& \\
3 \& \& 0 \arrow[ul] \arrow[dl] \\
\& 2 \arrow[ul] \&
\end{tikzcd}
\]
is the inclusion of vertex 3 in $\tilde{\Psi}$. If $\Ho_W(\Psi)$ is interpreted, in accordance with lemma 1, as $\mathcal{F}W_{\mathcal{F}}^{-1}$, this composite functor is none other than
\[
\begin{tikzcd}[ampersand replacement=\&]
X_1 \arrow[d] \& \\
S \& X_2 \arrow[l]
\end{tikzcd}
\qquad\longmapsto\qquad X_1\times_S X_2
\]
\printedpage{166}
% head: A. Grothendieck : Dérivateurs, Ch. XII, Édition des Universités de Montpellier II et Paris VII
% foot: Version 30 mars 2003
(for $\left(\begin{tikzcd}[ampersand replacement=\&] X_1 \arrow[d, "p_1"'] \& \\ S \& X_2 \arrow[l, "p_2"] \end{tikzcd}\right)\in \mathrm{Ob}\,\mathcal{F}$, i.e. $p_1,p_2\in\Fib_W$.)

We want to interpret this functor as a $W$-homotopic $\varprojlim$, i.e. viewed as a ``$W$-homotopic fiber product''. To be precise, we have this:

\msp{page 256}
\textbf{Corollary 1.} \textit{Consider $p_\Psi : \Psi\longrightarrow e$, and the functor}
\[
p_\Psi^* : \Ho_W(e)=\Ho_W\longrightarrow\Ho_W(\Psi),
\]
\textit{deduced by passage to the categories of fractions from the functor}
\[
\begin{array}{ccc}
\mathcal{M} & \longrightarrow & \mathcal{M}(\Psi) \\
X & \longrightarrow & \left(
\begin{tikzcd}[ampersand replacement=\&]
X \arrow[d] \& \\
X \& X \arrow[l]
\end{tikzcd}
\right).
\end{array}
\]
\textit{This functor admits the preceding functor $\varepsilon_3^*i_*$ as a right adjoint, i.e.}
\[
p_{\Psi\,*}=\varepsilon_3^*i_* : \Ho_W(\Psi)\longrightarrow\Ho_W.
\]

\textsc{Proof}. We can regard $p_\psi$ as the composite % [sic] p_\psi -> p_\Psi, the subscript is the capital \Psi everywhere else
\[
\Psi\xrightarrow{\ i\ }\tilde{\Psi}\xrightarrow{\ p_{\tilde{\Psi}}\ }e,
\]
and we therefore have
\[
p_\Psi^*\simeq i^*(p_{\tilde{\Psi}})^*.
\]
Thus the existence and the determination of a right adjoint for $p_\Psi^*$ will follow from the existence and the determination of right adjoints $i_*$, $(p_{\tilde{\Psi}})_*$ of $i^*$, $(p_{\tilde{\Psi}})^*$, via
\[
(p_\Psi)_*\simeq(p_{\tilde{\Psi}})_*i_*.
\]

\msp{page 257}
Thus corollary 1 follows from the formula
\[
p_{\tilde{\Psi}\,*}=\varepsilon_3^*,
\]
which itself is a special case of the

\printedpage{167}
% head: A. Grothendieck : Dérivateurs, Ch. XII, Édition des Universités de Montpellier II et Paris VII
% foot: Version 30 mars 2003

\textbf{Lemma.} \textit{Let $I$ be in} Cat\textit{, with final object $e$, consider the inclusion}
\[
\varepsilon : e\longrightarrow I,\qquad \textit{and } p : I\longrightarrow e.
\]
\textit{Let $(\mathcal{M},W)$ be a category of models, and consider}
\[
\Ho_W(I)\xleftarrow{\ p^*\ }\Ho_W.
\]
\textit{This functor admits a left adjoint} \texttt{[rather right]} \textit{$p_*$, which is isomorphic to $\varepsilon^*$.}

\textsc{Proof}. Consider the pair of adjoint functors
\[
\begin{tikzcd}[ampersand replacement=\&, column sep=large]
\mathcal{M}(I) \arrow[r, shift right, "p_*=\varprojlim\limits_I"'] \& \mathcal{M} \arrow[l, shift right, "p^*"']
\end{tikzcd}\ ;
\]
as $I$ admits the final object $e$, the functor $p_*$ is isomorphic to $\varepsilon^*$. These functors $p^*$, $\varepsilon^*$ are compatible with the localizers $W(I)$, $W$, so we conclude by the lemma on page 248.

It only remains to prove lemmas 1 and 2. I will do it tomorrow.

\msp{page 258}
Today, I still feel like giving the sorites of $W$-cartesian squares in a

\textbf{Proposition 19 :} \textit{Let $\begin{tikzcd}[ampersand replacement=\&] X_1 \arrow[d, "p_1"'] \& X_3 \arrow[l, "f_1"'] \arrow[d, "f_2"] \\ S \& X_2 \arrow[l, "p_2"] \end{tikzcd}$ be a commutative square in $\mathcal{M}$. Consider factorizations of $p_1$, $p_2$ as}
\[
X_1\xrightarrow{\ i_1\ }\bar{X}_1\xrightarrow{\ \bar{p}_1\ }S,\qquad X_2\xrightarrow{\ i_2\ }\bar{X}_2\xrightarrow{\ \bar{p}_2\ }S,
\]
\textit{$i_1,i_2\in W$, $\bar{p}_1,\bar{p}_2\in F_W$} \texttt{[$=$ Fib$_W$]}\textit{, and consider the arrows induced by $X_3\longrightarrow X_1\times_S X_2$ (deduced from $Q$) via $\begin{tikzcd}[ampersand replacement=\&, row sep=small] \& \bar{X}_1\times_S X_2 \\ X_1\times_S X_2 \arrow[ur] \arrow[r] \arrow[dr] \& \bar{X}_1\times_S\bar{X}_2 \\ \& X_1\times_S\bar{X}_2 \end{tikzcd}$, so that we find a commutative diagram}
\[
\begin{tikzcd}[ampersand replacement=\&, row sep=small]
\& \bar{X}_1\times_S X_2 \arrow[d, "j_2"'] \\
X_3 \arrow[ur, "\pi_1"] \arrow[r, "\pi_0"] \arrow[dr, "\pi_2"] \& \bar{X}_1\times_S\bar{X}_2 \\
\& X_1\times_S\bar{X}_2, \arrow[u, "j_1"]
\end{tikzcd}
\]
\textit{where $j_1=i_1\times_S\mathrm{id}_{\bar{X}_2}$ and $j_2=\mathrm{id}_{\bar{X}_1}\times i_2$ are in $W$ (since $i_1$ and $i_2$ are, and $\bar{X}_1$ and $\bar{X}_2$ are in $\underline{\mathrm{Fib}}_W S$). We therefore have}
\[
(\pi_1\in W)\qquad\iff\qquad(\pi_2\in W)\qquad\qquad(\iff\ \pi_0\in W),
\]
\textit{which proves that the condition $\pi_1\in W$, which seems to depend on the choice of the factorization through $i_1$, does not depend on it, and likewise for $\pi_2\in W$, and for $\pi_0\in W$ (which seems to depend on both $i_1$, and $i_2$).}

\printedpage{168}
% head: A. Grothendieck : Dérivateurs, Ch. XII, Édition des Universités de Montpellier II et Paris VII
% foot: Version 30 mars 2003

\msp{page 259}
\textbf{Definition.} \textit{When these conditions are satisfied, we say that the square $Q$ is $W$-cartesian.}

\textbf{Corollary 1.} \textit{If $Q$ is $W$-cartesian, then the image of $Q$ in} $\Ho_W(\tilde{\Psi})$ \textit{is in the essential image of}
\[
i_* : \Ho_W(\Psi)\longrightarrow\Ho_W(\tilde{\Psi}),
\]
\textit{i.e. the canonical arrow in} $\Ho_W(\tilde{\Psi})$
\[
Q\longrightarrow i_*i^*Q
\]
\textit{is} \texttt{[an]} \textit{isomorphism. The converse is true if $W$ is strongly saturated.}

\textbf{Corollary 2.} \textit{Consider a commutative diagram in $\mathcal{M}$}
\[
\begin{tikzcd}[ampersand replacement=\&]
X_0' \arrow[d] \& X_1' \arrow[l] \arrow[d] \arrow[dl, phantom, "\textit{1}"] \& X_2' \arrow[l] \arrow[d] \arrow[dl, phantom, "\textit{2}"] \\
X_0 \& X_1 \arrow[l] \& X_2, \arrow[l]
\end{tikzcd}
\tag*{(D)}
\]
\textit{whence two commutative squares 1, 2, and a composite square}
\[
\begin{tikzcd}[ampersand replacement=\&]
X_0' \arrow[d] \& X_2' \arrow[l] \arrow[d] \arrow[dl, phantom, "\textit{3}"] \\
X_0 \& X_2. \arrow[l]
\end{tikzcd}
\]
\textit{Suppose square 1 is $W$-cartesian. Then square 2 is $W$-cartesian if and only if square 3 is.}

\textsc{Proof}. Consider a factorization of $\begin{tikzcd}[ampersand replacement=\&, row sep=small] X_0' \arrow[d] \\ X_0 \end{tikzcd}$ as
\[
\begin{tikzcd}[ampersand replacement=\&, column sep=small, row sep=small]
\bar{X}_0' \arrow[dr] \& X_0' \arrow[l] \arrow[d] \\
\& X_0
\end{tikzcd}
\]
with $(X_0'\longrightarrow\bar{X}_0')\in W$, and $(\bar{X}_0'\longrightarrow X_0)\in\Fib_W$. We can then replace $X_0'$ by $\bar{X}_0'$ in (D), obtaining a diagram $(\bar{D})$, and we see at once
\printedpage{169}
% head: A. Grothendieck : Dérivateurs, Ch. XII, Édition des Universités de Montpellier II et Paris VII
% foot: Version 30 mars 2003
\msp{page 260}
that each of the squares 1(D), 2(D), 3(D) is $W$-cartesian if and only if the corresponding square $1(\bar{D})$, $2(\bar{D})$ or $3(\bar{D})$ is. This shows that we can assume $X_0\longrightarrow S$ to be in $\Fib_W S$. Consider then the commutative diagram deduced from (D)
\[
\begin{tikzcd}[ampersand replacement=\&, column sep=small]
X_0' \arrow[d] \& \tilde{X}_1'\overset{\text{def}}{=}X_0'\times_{X_0}X_1 \arrow[l] \arrow[d] \& X_1' \arrow[l, "\alpha"'] \& X_2' \arrow[l] \arrow[d] \\
X_0 \& X_1 \arrow[l] \& \& X_2, \arrow[ll]
\end{tikzcd}
\]
where $\alpha\in W$ by virtue of the $W$-cartesian hypothesis on square 1. We see again that square 2 is $W$-cartesian if and only if the corresponding square with $X_1'$ replaced by $\tilde{X}_1'$ is, as for squares 1, 3, they do not change under this modification. This shows that we can assume square 1 \textit{cartesian.} Consider then the diagram
\[
\begin{tikzcd}[ampersand replacement=\&]
X_0' \arrow[d, "\mathrm{Fib}_W"'] \& X_1' \arrow[l] \arrow[d, "\mathrm{Fib}_W"] \arrow[dl, phantom, "\text{cart.}"] \& \tilde{X}_2'\overset{\text{def}}{=}X_1'\times_{X_1}X_2 \arrow[l] \arrow[d, "\mathrm{Fib}_W"] \arrow[dl, phantom, "\text{cart.}"] \& X_2' \arrow[l, "\beta"'] \\
X_0 \& X_1 \arrow[l] \& X_2. \arrow[l]
\end{tikzcd}
\]
Then proposition 19 and the definition show that square 2 is $W$-cartesian if and only if $\beta\in W$, and likewise for square 3. Which completes the proof. % [sic] seulenment -> seulement, misspelling

\textbf{Corollary 2.} \textit{Let $\begin{tikzcd}[ampersand replacement=\&] X_1 \arrow[d, "p_1"'] \& X_3 \arrow[l, "f_2"'] \arrow[d, "f_1"] \\ S \& X_2 \arrow[l, "p_2"] \end{tikzcd}$ be a $W$-cartesian square. Then}
\[
\begin{aligned}
p_1\in W\quad&\Longrightarrow\quad f_1\in W\\
p_2\in W\quad&\Longrightarrow\quad f_2\in W,
\end{aligned}
\]
\textit{more precisely, if $p_i\in W$, then $Q$ is $W$-cartesian if and only if $f_i\in W$.}

\msp{page 261}
Suppose $p_1\in W$. Then in the factorization of $p_1$ as $X_1\longrightarrow\bar{X}_1\longrightarrow S$, we can take $\bar{X}_1=S$. Thus the square is $W$-cartesian if and only if
\[
(X_3\longrightarrow\bar{X}_1\times_S X_2=X_2)\in W,
\]
i.e. $f_1\in W$, q.e.d.

Of course, we cannot in general assert the reverse implications: $f_i\in W$ by no means implies $p_i\in W$.

\end{document}
